用 \texttt{sps} 命名空间，按 GNU C++23 编译。数组长 \(m=2^n\)（\(0\le n\le30\)），
\texttt{f[s]} 表示掩码 \(s\) 的系数，\texttt{f[0]} 是空集项。二元运算的两数组等长。
模数 \(P\) 为满足 \(\max(n,2)<P<2^{30}\) 的素数，系数取 \([0,P)\)。
不交子集卷积的单位元仅空集项为 \(1\)，其余项均为 \(0\)。

\paragraph*{子集卷积为什么要分层}
子集和的点乘会同时计入有交和无交的集合对。对 \(f,g\) 分别按集合大小求子集和，
得到 \(A_i[S],B_j[S]\)。固定 \(S\)，按次数相乘得到 \(C_k[S]\)。
\[
    A_i[S]=\sum_{\substack{U\subseteq S\\|U|=i}} f[U],
    \qquad
    C_k[S]=\sum_{i=0}^{k} A_i[S]B_{k-i}[S].
\]
乘积 \(f[U]g[V]\) 落在第 \(k=|U|+|V|\) 层，并对每个 \(S\supseteq U\cup V\) 贡献一次。
逐层做子集反演后，这一项只留在 \(S=U\cup V\)。此时取 \(k=|S|\)，恰好筛出
\(|U|+|V|=|U\cup V|\) 的项，即 \(U\cap V=\varnothing\)。
因此应先反演，再取对应层。\texttt{mul} 倒序计算 \(k\) 以复用 \(A\)，时间 \(O(n^2 2^n)\)，空间 \(O(n2^n)\)。

\paragraph*{\texttt{inv} 与 \texttt{exp} 怎样拆最高位}
设旧数组长 \(h\)，加入新元素 \(x\) 后把输入写成 \(F=A+xB\)。
对 \(0\le s<h\)，低半段 \(A[s]=f[s]\)，高半段 \(B[s]=f[h+s]\)，后者已去掉新位。
以下乘法为不交子集卷积。两部分若都含 \(x\) 就会相交，故 \(x^2=0\)。

求逆时，已知低半段 \(G=A^{-1}\)，设答案为 \(G+xH\)。展开
\((A+xB)(G+xH)=1+x(AH+BG)\)，令一次项为零，得 \(H=-GBG\)。
代码从 \(f[0]^{-1}\) 开始，每轮把 \(-GBG\) 接到 \(G\) 后面。

求指数时，已知低半段 \(G=\exp A\)。由 \(x^2=0\) 得 \(\exp(xB)=1+xB\)，
故 \(\exp F=G(1+xB)=G+xGB\)。代码从 \(1\) 开始，每轮只需接上 \(GB\)。
两种运算的卷积规模逐轮翻倍，总时间均为 \(O(n^2 2^n)\)。
